\chapter{Grammar and glamour}


\subsection{Even it's never approved}

In 2012, Cory Doctorow wrote on \textit{Boing Boing}, ``But if they test it on kids -- even it's never approved for use on kids! -- the folks at the FDA will extend their patent by another six months.''

My English-language students say and write things like \textit{even it's never approved} (meaning `even if it's never approved') all the time, but it's ungrammatical for me -- likely for you too -- and I correct them. Doctorow's post was the first time I'd seen it from a proficient speaker of English, so I emailed him to ask if it was intentional (not to correct him). ``Yup,'' he responded. ``It's colloquial, but expressive, unambiguous, and colorful.''\footnote{Personal communication, July 6, 2012.}

There's an \textit{even} + pronoun construction that does work, one with stress on the pronoun, as in \textit{even \textbf{I} know that's ungrammatical.} But that's not happening here. In this innovation, \textit{even} takes on a conditional meaning and becomes a preposition that takes a clausal complement. \textit{If} is missing, but the sense and syntax is otherwise unchanged.

You may accuse me of arbitrary discrimination against English-language learners, but if you did (even you did?) I would protest. If Doctorow had written \textit{even if it's never approve}, I would have instantly recognized it as a typo. Same if he'd said, \textit{even if he's never approved}. Same for a million other mistakes. It's simply not the case that a proficient speaker of a language can get away with anything.

A grammatical innovation will only work if it achieves some kind of semantic or pragmatic goal, something different from what the established form achieves, an effect intended by the innovator. That could be humour, group membership, irony, emphasis, (in)formality, (in)directness, intimacy, immediacy, skepticism, or any number of other effects. And to do that, there needs to be some establish connection between the innovation and the effect, some linguistic or cultural point of reference.

A grammatical innovation that fails to do these things is ungrammatical.

\subsection{Tracksuits}

By the year 2000, the tracksuit was clearly unfashionable, a relic of a past era of bold and sporty looks that no longer resonated with the sleek, minimalist trends of the new millennium. Since the start of the pandemic, though, the matching sweatpants-and-sweatshirt look has gained currency \citep{friedman2024}.

\begin{quote}
    As with most things, the details matter. Wear a sweatsuit with a sweatband and sneakers and you look like a victim of the endless cycle of New Year's workout resolutions. Wear it with a nice handbag and a trench coat and you look like a modern Jean Harlow out to take the air.\hfill(Vanessa Friedman)
\end{quote}

Just as a fashion influencer can take this once-derided outfit and transform it into a statement of effortless chic, a linguistically savvy individual can take a phrase that might be considered a grammatical misstep and turn it into a mark of stylistic flair. The phrase \textit{a pair glasses}, when uttered by someone grappling with the intricacies of English, will likely be dismissed as ungrammatical. Yet, when used by a proficient speaker with the linguistic sense to use it intentionally to achieve some goal, it becomes a deliberate stylistic choice, an intentional innovation.

\textit{I have a pair glasses}, for instance, might signal irony in reply to ``I have a dozen tracksuits,'' \textit{dozen} establishing the template on which the new \textit{pair} construction is built.\footnote{Most collectives, like \textit{group}, \textit{bunch}, \textit{lot}, \textit{pile}, \textit{heap}, and \textit{shitload} license an \textit{of} preposition phrase complement (e.g., \textit{a group of words}). \textit{Dozen} is different: *\textit{a dozen of words} is ungrammatical, though \textit{a dozen of them} is fine. In this way, \textit{a dozen} is like the determinatives \textit{a little} and \textit{a few}, and \textit{a dozen of them} is partitive, where the object of the \textit{of} preposition phrase has to be definite.}

The tracksuit's journey from the back of the wardrobe to the forefront of fashion showcases how context and the wearer's choices can completely alter the perception of an item. When paired with the right accessories and worn with purpose, the once unfashionable item becomes glamorous, the once ungrammatical expression becomes the hot new slang.

\subsection{The glory of \textit{enshittification}}

\textit{Enshittification} was the American Dialect Society 2023 selection for Word of the Year \citep{Roberts2024}, and Cory Doctorow is credited with popularizing it. In early 2023, he wrote,

%Jan 5, 2024 Vote took place at the Sheraton New York Times Square hotel.

\begin{quote}
    Here is how platforms die: first, they are good to their users; then they abuse their users to make things better for their business customers; finally, they abuse those business customers to claw back all the value for themselves. Then, they die. I call this \textit{enshittification},\hfill\citep{doctorow2023}
\end{quote}

It's a perfectly cromulent word, the kind any linguist loves: new, creative, flexible, and evocative. When it came to the attention of linguist Nancy Friedman, she said it was ``primed for neologistic success'' \citep{friedman2024sweary}. Ben Zimmer wrote, ``From the time that it first appeared in Doctorow's posts and articles, the word had all the markings of a successful neologism, being instantly memorable and adaptable to a variety of contexts.''

Part of what makes \textit{enshittification} so appealing is its anchoring in a linguistic and cultural history. Morphologically, it starts with an everyday term and then accessorizes it with a gaudy array of familiar affixes, each elevating its lexical glory.

The base word here is \textit{shit}, a noun with Old English roots, specifically \textit{scitte}, meaning `diarrhea', which is itself derived from the Proto-Germanic \textit{skit--}. Needless to say, over centuries, \textit{shit} has evolved in English as a vulgar term for feces, but also as a versatile word used in a variety of idiomatic expressions and colloquialisms.

To this indelicate base, he adds \textit{--ify}, from the Latin \textit{--ficare}, which in turn is derived from \textit{facere}, meaning `make' or `do'. This suffix turns nouns or adjectives into verbs, denoting the process of making something have a particular quality or be in a particular state. Doctorow could have simply said the platforms \textit{shittify} themselves, but its unlikely that the word would have had the appeal of \textit{enshittification}. It's a little too literal, a little too sharp, a little too base. It needs more decoration.

\textit{--ic}, an adjective-forming suffix denoting a characteristic or quality, attaches readily to \textit{--ify}, as in \textit{terrific}, but an adjective won't do. So on he goes.

% TODO: verify that the "Three suffixes is grand..." quote is genuinely from Doctorow — flagged by review-board as possibly invented; if it isn't sourced, recast as Brett's own gloss rather than a Doctorow quotation.
By the curious logic of English, \textit{--ific} attracts the suffix \textit{--ate}, resulting in \textit{shittificate}, the opposite of \textit{beatificate}, re-verbing the word. ``Three suffixes is grand,'' says Doctorow, ``but I can do more. \textit{Shittificate} has three, but I can do four'' (with apologies to Dr. Seuss).

\textit{--tion} is a noun-forming suffix that turns verbs into nouns. When attached to \textit{--ificate}, as in \textit{magnification}, it creates a noun that indicates the action or process of making something have a particular quality. He's now gone full circle from the noun \textit{shit} to the noun \textit{shittification}, and \textit{shittification} could denote what he's after. Doctorow, though, is still not done.

In a bold step, audacious even, he adds the completely redundant \textit{en--} prefix. This prefix, typically used to turn nouns into verbs (as in \textit{enlighten} from \textit{light}), here seems superfluous. Yet, its inclusion is a masterstroke in neologism, enhancing the word's dramatic flair. \textit{Enshittification} is not just about becoming shit or the act of making something shitty -- it's about the willful and relentless driving of a platform into the ground by stuffing it with shit.

Doctorow uses the familiar to produce the sublime. But it's not just the use of familiar affixes but the particular affixes used that have resonance with other successful neologisms. \textit{Embiggen} appeared on ``The Simpsons'' in 1996 and stuck around, finding a spot in the \textit{Oxford English dictionary} in 2017 and in the \textit{Merriam-Webster dictionary} in 2018. Doctorow was almost certainly aware of this when he chose the \textit{en--}/\textit{em--} prefix.

Language aside, \textit{enshittification} perfectly captures the tech world's slide from revolutionary to reviled. We've all experienced that moment when a platform, once a beacon of digital innovation, becomes just crummy, betraying its original ethos, turning into something users hardly recognize. We've all experienced it. Doctorow crystallized it.

He's a blogger, a journalist, a science fiction author, an expert language user, and somebody really tapped into the social zeitgeist. He knows what he's doing. The final form, \textit{enshittification}, thus becomes a linguistic ornament, an embellished word capturing a complex social phenomenon in a single, extravagant term. And all of that skill and intentionality lead to the acceptance and recognition that the word has received.

